% Cover letter for CEI VinUni - Research Associate, Environmental Vulnerability & Governance

\documentclass[10pt]{article}
\usepackage[T1]{fontenc}
\usepackage[utf8]{inputenc}
\usepackage{libertine}
\usepackage[margin=1.7cm]{geometry}
\usepackage[hidelinks]{hyperref}
\usepackage{xcolor}

\setlength{\parindent}{0pt}
\setlength{\parskip}{0.58em}
\pagestyle{empty}

\begin{document}

\begin{flushright}
Huy Le Vu\\
\href{mailto:huymaeco@gmail.com}{huymaeco@gmail.com}\\
\href{https://huylvu.github.io}{huylvu.github.io}\\
June 3, 2026
\end{flushright}

Selection Committee\\
Center for Environmental Intelligence\\
VinUniversity

Dear Selection Committee,

I am writing to apply for the Research Associate position on Environmental Economics, Public Policy, and Environmental Hazards at the Center for Environmental Intelligence, VinUniversity. The project is a strong fit for my background in applied economics and policy research in Viet Nam, especially household and firm survey preparation, fieldwork-facing data systems, environmental and climate-related analysis, Vietnamese policy context, and academic writing.

I do not yet hold a Master's degree, and I want to state that clearly. At the same time, I believe my experience makes me a qualified candidate for this role. From March 2024 to May 2026, I worked as a Research Assistant at the Development and Policies Research Center (DEPOCEN), supporting projects for UNDP, the World Bank, ADB, USAID/CRS, Economics for Health, and other partners. I also have two peer-reviewed journal publications, a strong undergraduate record in International Economics at Foreign Trade University, and practical experience turning research questions into survey instruments, analytical datasets, tables, figures, reports, and manuscript-facing evidence.

Several parts of my experience map directly onto your project. I led the Viet Nam Multi-Sector Assessment of the Typhoon Yagi Macroeconomic Impact for UNDP in 2024, which strengthened my understanding of climate shocks, regional exposure, policy response, and evidence needs after a major disaster. I also supported KOICA Green Circular Economy baseline and inception work, including sampling logic, household questionnaire structure, livelihood screening, skip-pattern review, and survey documentation. In the CRS/USAID Disability Project, I worked across repeated survey rounds, supporting data preparation, quality checks, and research documentation. These projects made me comfortable with the practical side of research: sampling frames, instrument logic, field constraints, data quality, and transparent documentation of methodological choices.

I would also bring a strong fit for the analytical and writing responsibilities of the role. My published articles in the \textit{Journal of Economics and Development} and \textit{Thailand and The World Economy} reflect experience with empirical analysis, academic writing, revision, formatting, and publication workflows. I am proficient in Stata, R, and QGIS, and I also use Python, \LaTeX, Excel, and Git for reproducible research workflows. I have worked with geospatial and environmental data, including shapefiles, administrative boundaries, temperature, rainfall, night-time lights, CO2, greenhouse gas emissions, and PM2.5 indicators. I am fluent in Vietnamese and comfortable working with Vietnamese policy materials, which is essential for policy-document analysis and field coordination.

The role's combination of household survey implementation, Community Vulnerability Indicators and Household Vulnerability Index work, Vietnamese environmental governance, and manuscript development is exactly the kind of research environment where I can contribute immediately while continuing to grow. I would be ready to support the second-phase household survey, integrate new data with the existing dataset, assist with index construction and sensitivity checks, synthesize Vietnamese and international literature, and help prepare academic outputs for journal and conference submission.

Thank you for considering my application. I would welcome the opportunity to discuss how my research experience, publication record, and fieldwork-facing data skills can contribute to the Environmental Vulnerability and Governance project at CEI.

Sincerely,\\
Huy Le Vu

\end{document}
